\subsection{Criterios de selección}

Para integrar compresión de datos en TMFQSfullstate (simulador de estado completo de circuitos
cuánticos), se definen los siguientes criterios técnicos, ordenados por prioridad:

\begin{table}[H]
    \centering
    \vspace{-2mm}
    \caption{\\Criterios para la selección de compresores sin pérdida en TMFQSfullstate}
    \label{tab:criterios_compresion}
    \begin{tabular}{lm{4.1cm}m{0.62\linewidth}c}
        \hline
        \multicolumn{1}{l}{} &
        \multicolumn{1}{c}{\textbf{Criterio}} &
        \multicolumn{1}{c}{\textbf{Explicación}} &
        \multicolumn{1}{c}{\textbf{Peso}} \\ \hline

        C1 & \centering Ratio de compresión &
        Un mayor factor de compresión permite representar el vector de estado con menor memoria,
        habilitando la simulación de más qubits con los mismos recursos. En simulación de estado
        completo, la memoria suele ser el cuello de botella principal, por lo que maximizar la
        reducción de memoria sin pérdida es el objetivo dominante; este énfasis es consistente con
        la motivación de compresión en simulación cuántica a gran escala reportada por Wu y cols.
        (2019b). &
        0.40 \\

        C2 & \centering Velocidad de compresión y descompresión &
        La simulación puede requerir ciclos frecuentes de compresión/descompresión, por lo que el
        algoritmo debe minimizar latencia y maximizar throughput. Se priorizan compresores con
        throughput de centenas de MB/s (o más) y con descompresión aún más rápida para no retrasar
        la recuperación de datos (lz4 contributors, s.f.; facebook/zstd contributors, s.f.; Snappy
        project, s.f.; MaskRay, 2025). &
        0.20 \\

        C3 & \centering Bajo overhead de memoria &
        Dado el crecimiento exponencial del estado, el compresor no debe demandar estructuras
        internas (buffers/diccionarios/tablas) que incrementen sustancialmente el uso total de RAM.
        Se favorecen métodos que operen en streaming o por bloques pequeños (lz4 contributors, s.f.;
        Zstandard project, s.f.). &
        0.15 \\

        C4 & \centering Compatibilidad con procesamiento paralelo &
        En escenarios multi-core o distribuidos (p. ej., MPI en HPC), interesa que el compresor
        soporte multihilo de forma nativa o, como mínimo, que sea fácilmente paralelizable por
        bloques sin degradaciones significativas de ratio o throughput (facebook/zstd contributors, s.f.). &
        0.15 \\

        C5 & \centering Integración en C++ con TMFQSfullstate &
        El algoritmo debe estar disponible como librería utilizable desde C/C++ y ser suficientemente
        madura, con soporte y licencia permisiva, para minimizar riesgo de integración y mantenimiento
        (lz4 contributors, s.f.; facebook/zstd contributors, s.f.; google/snappy contributors, s.f.). &
        0.10 \\

        \hline
    \end{tabular}
    \vspace{-2mm}
\end{table}

Los pesos se asignan para operacionalizar el orden de prioridad definido: se privilegia el ratio
(C1) como objetivo dominante, seguido por el costo temporal (C2), y luego por restricciones de
ingeniería relevantes para ejecución a gran escala e integración (C3--C5).
